\documentclass[10pt,a4paper]{amsart}
\usepackage[margin=19mm]{geometry}
\usepackage{amsmath,amssymb,mathtools}
\usepackage[scr=rsfs,cal=euler]{mathalfa}
\usepackage{xfrac}
\usepackage[unicode,hidelinks,bookmarks=false]{hyperref}
\newcommand{\ie}{i.e.\ }
\newcommand{\cf}{cf.\ }
\newcommand{\resp}{respectively\ }
\newcommand{\Subsection}{\subsection}
\providecommand{\nobreakdash}{\nobreak\hbox{-}}
\setlength{\emergencystretch}{2em}
\multlinegap0pt
\usepackage{amssymb}
\usepackage{float}
\usepackage[shortlabels]{enumitem}
\newtheorem{theorem}{Theorem}
\title{Vertex evaluation of multiplex graphs using Forman Curvature}
\author{Taiki Yamada}
\date{Source: arXiv:2504.17286v2 (CC BY 4.0)}
\begin{document}
\maketitle

\noindent\emph{Source and licence.} Vertex evaluation of multiplex graphs using Forman Curvature. Version arXiv:2504.17286v2 (CC BY 4.0), distributed by arXiv under the Creative Commons Attribution 4.0 International licence, http://creativecommons.org/licenses/by/4.0/, as recorded in the arXiv licence metadata for this exact version and shown on the abstract page https://arxiv.org/abs/2504.17286v2. Full licence notice: \texttt{licenses/arxiv-2504.17286-CC-BY-4.0.txt}.\par\smallskip

\noindent\textit{Editorial scope.} Counted unit: the article's theorem lower (source0.tex lines 449--461). The article's complete original proof of it is delivered with it (source0.tex lines 463--486, one \texttt{proof} environment of 24 lines). No mathematics is written, repaired or re-derived: the statement and the proof are the article's own lines, in the article's own notation, and no retained line is substituted or re-flowed.  No further source-local context is needed: every cross-reference of every retained line resolves inside the two delivered spans.  Nothing was omitted from any retained span, so the package carries no omission record.  No retained line contains a citation, so the wrapper carries no bibliography and no bibliography entry.  No retained line contains a figure environment, an image-inclusion command or drawing code, so no image is delivered and the package declares no asset dependency.  Editorial formatting only, in addition: an AMS \texttt{amsart} preamble that restates the article's own declarations and macros used by the retained lines, namely the environments \texttt{theorem} on separate counters, together with the standard packages those lines need.  The retained environments print in this document as Theorem 1 in order of appearance, whereas the article prints the same items with the numbering of its own full document.  The complete native source of the article as distributed in the arXiv e-print for this exact version is bundled unchanged as \texttt{sources/177014/source0.tex}.
\begin{theorem}
\label{lower}
    Let $CG = (V_{CG}, E_{CG}, V, \mathcal{L}, w, m)$ be a compile graph. 
    For any inter-layer edge $e =(x^i, x^j) \in E_C$ with $W(x^i) \leq W(x^j)$, we then have
    \begin{eqnarray}
    \label{Formanvalue}
        F((x^i,x^j)) 
        &\geq& -(m(x^i)+m(x^j)) |\Gamma^{i,j}_{C,-} (x)|\max \left\{ \cfrac{W(x^i)}{W(x^l)} \mid x^l \in \Gamma^{i,j}_{C,-}(x)\right\} \nonumber\\
        &\ & - m(x^j)\left(\cfrac{W(x^i)}{W(x^j)} -1 \right)\\
        &\geq& -(m(x^i)+m(x^j)) |\Gamma^{i,j}_{C,-} (x)|\max \left\{ \cfrac{W(x^i)}{W(x^l)} \mid x^l \in \Gamma^{i,j}_{C,-} (x)\right\}.
    \end{eqnarray}
    Equality holds in the first inequality line if and only if $\Gamma^{i,j}_C (x)= \Gamma^{i,j}_{C,-}(x)$. Equality holds in the second inequality line if and only if $W(x^i) = W(x^j)$.
\end{theorem}

\begin{proof}
    Note that for any subset $A$ of $X$, if the number of elements in $A$ decreases by one, the number of elements in $X\setminus A$ increases by one.

    If there exists a vertex $x^k \in V_{CG}$ such that $x^k \notin \Gamma_{C,-}^{i,j}(x)$, we consider another compile graph $\bar{CG}$ where only the value of $W(x^k)$ is modified to ensure that $x^k$ is included in $\Gamma_{C,-}^{i,j}(x)$. Let the modified value of $W(x^k)$ be $\bar{W}(x^k)$. In this case, the difference in the Forman curvature is calculated as follows:
    \begin{eqnarray*}
        F^{CG}((x^i,x^j)) - F^{\bar{CG}}((x^i,x^j)) = m(x^i) \left( \cfrac{W(x^i)}{\bar{W}(x^k)} - 1\right) \geq 0,
    \end{eqnarray*}
    which implies $F^{CG}((x^i,x^j)) \geq F^{\bar{CG}}((x^i,x^j))$. Thus, the larger the cardinality of $\Gamma_{C,-}^{i,j}(x)$, the smaller is the Forman curvature.

    Conversely, if there exists a vertex $x^k \in V_{CG}$ such that $x^k \notin \Gamma_{C,+}^{i,j}(x)$, we consider another compile graph $\hat{CG}$ where only the value of $W(x^k)$ is modified to ensure that $x^k$ is included in $\Gamma_{C,+}^{i,j}(x)$. Let the modified value of $W(x^k)$ be $\hat{W}(x^k)$. In this case, the difference in Forman curvature is calculated as follows:
    \begin{eqnarray*}
        F^{CG}((x^i,x^j)) - F^{\hat{CG}}((x^i,x^j)) = m(x^j) \left( \cfrac{W(x^i)}{W(x^j)} - \cfrac{W(x^i)}{W(x^k)}\right) \leq 0,
    \end{eqnarray*}
    which implies $F^{CG}((x^i,x^j)) \leq F^{\hat{CG}}((x^i,x^j))$. Thus, the larger the cardinality of $\Gamma_{C,+}^{i,j}(x)$, the larger  the Forman curvature.

    Based on these results, the minimum value of the Forman curvature occurs when the cardinality of $\Gamma_{C,-}^{i,j}(x)$ coincides with that of $\Gamma_C^{i,j}(x)$. The minimum value is calculated as follows:
    \begin{eqnarray*}
        F((x^i,x^j)) 
        &\geq& -(m(x^i)+m(x^j)) |\Gamma^{i,j}_{C,-} (x)|\max \left\{ \cfrac{W(x^i)}{W(x^l)} \mid x^l \in \Gamma^{i,j}_{C,-}(x)\right\} \nonumber\\
        &\ & - m(x^j)\left(\cfrac{W(x^i)}{W(x^j)} -1 \right).
    \end{eqnarray*}

    Next, we examine the maximum value of $W(x^i)/W(x^j)$. Given the condition that $W(x^i) \leq W(x^j)$, the maximum value of $W(x^i)/W(x^j)$ is 1 when $W(x^i)=W(x^j)$. Thus, the lower bound of the Forman curvature was proven.
\end{proof}

\end{document}
